\documentclass[12pt,a4paper]{article}

\usepackage[T1]{fontenc}
\usepackage[utf8]{inputenc}
\usepackage[brazilian]{babel}
\usepackage{times}
\usepackage{geometry}
\usepackage{setspace}
\usepackage{indentfirst}
\usepackage{titlesec}
\usepackage{fancyhdr}
\usepackage{microtype}
\usepackage[hidelinks]{hyperref}

% Configuração baseada na ABNT NBR 14724:2024 e no Manual da UNIFAL-MG.
\geometry{top=3cm,left=3cm,right=2cm,bottom=2cm,headheight=14.5pt}
\onehalfspacing
\setlength{\parindent}{1.25cm}
\setlength{\parskip}{0pt}
\setlength{\emergencystretch}{3em}

\titleformat{\section}
  {\normalfont\bfseries\normalsize}
  {\thesection}{0.5em}{\MakeUppercase}
\titlespacing*{\section}{0pt}{1.5\baselineskip}{1.5\baselineskip}

\pagestyle{fancy}
\fancyhf{}
\fancyhead[R]{\fontsize{10}{12}\selectfont\thepage}
\renewcommand{\headrulewidth}{0pt}

\begin{document}
\pagenumbering{gobble}
\hypersetup{pageanchor=false}

% CAPA
\begin{titlepage}
    \begin{center}
        \textbf{UNIVERSIDADE FEDERAL DE ALFENAS -- UNIFAL-MG}\\
        \textbf{BACHARELADO EM CIÊNCIA DA COMPUTAÇÃO}

        \vspace*{5cm}

        \textbf{JEANN VICTOR BATISTA}

        \vfill

        \textbf{RESUMO DO ARTIGO ``A NEW HUMAN-BASED METAHEURISTIC
        OPTIMIZATION METHOD BASED ON MIMICKING COOKING TRAINING''}

        \vfill

        \textbf{ALFENAS -- MG}\\
        \textbf{2026}
    \end{center}
\end{titlepage}

% FOLHA DE ROSTO
\begin{titlepage}
    \begin{center}
        \textbf{JEANN VICTOR BATISTA}

        \vspace*{5cm}

        \textbf{RESUMO DO ARTIGO ``A NEW HUMAN-BASED METAHEURISTIC
        OPTIMIZATION METHOD BASED ON MIMICKING COOKING TRAINING''}
    \end{center}

    \vspace*{3cm}

    \hfill
    \begin{minipage}{0.55\textwidth}
        \singlespacing
        Relatório apresentado à disciplina DCE770 -- Heurísticas e
        Metaheurísticas, do curso de Bacharelado em Ciência da Computação da
        Universidade Federal de Alfenas, como requisito da primeira etapa do
        Trabalho Final.

        \vspace{\baselineskip}

        Professor: Prof. Dr. Iago Augusto de Carvalho.
    \end{minipage}

    \vfill

    \begin{center}
        \textbf{ALFENAS -- MG}\\
        \textbf{2026}
    \end{center}
\end{titlepage}

% SUMÁRIO
\tableofcontents
\clearpage

% A capa não é contada; folha de rosto e sumário são contados, mas não numerados.
\hypersetup{pageanchor=true}
\pagenumbering{arabic}
\setcounter{page}{3}

\section{Identificação e contextualização do artigo}

O artigo \textit{A new human-based metaheuristic optimization method based on
mimicking cooking training}, de Eva Trojovská e Mohammad Dehghani, foi
publicado em 2022 no periódico de acesso aberto \textit{Scientific Reports}.
O trabalho apresenta o \textit{Chef-Based Optimization Algorithm} (CBOA), uma
meta-heurística populacional inspirada no processo de formação de cozinheiros
em cursos de culinária. O CBOA integra, portanto, a categoria das
meta-heurísticas baseadas em atividades humanas, ao lado de métodos como o
TLBO e o WSO. Sua finalidade é resolver problemas de otimização nos quais
métodos determinísticos podem ser inadequados em razão de não linearidade,
não convexidade, alta dimensionalidade, descontinuidade ou grande número de
ótimos locais (Trojovská; Dehghani, 2022).

\section{Problema estudado e motivações}

O problema geral considerado é encontrar, em um espaço de busca limitado, os
valores das variáveis que minimizam uma função objetivo. Em aplicações reais,
enumerar todas as soluções pode ser inviável e técnicas baseadas em gradientes
nem sempre podem ser empregadas. As meta-heurísticas utilizam busca estocástica
para produzir soluções próximas do ótimo global sem exigir derivadas. Contudo,
nenhum algoritmo é superior em todas as classes de problemas, conforme o
teorema \textit{No Free Lunch} (Wolpert; Macready, 1997). Essa ausência de um
otimizador universal constitui a motivação apresentada pelos autores para
desenvolver uma nova estratégia que equilibre exploração global e exploração
intensiva das regiões promissoras.

\section{Contribuições apresentadas}

A contribuição central é a formulação matemática do CBOA, que associa cada
solução candidata às habilidades de uma pessoa em treinamento culinário. A
população é reordenada a cada iteração e dividida entre chefs instrutores, que
correspondem aos indivíduos mais bem avaliados, e estudantes. Os autores
formalizam duas estratégias para atualizar os instrutores e três estratégias
para atualizar os estudantes. Também apresentam fluxograma, pseudocódigo,
análise de complexidade, análise de sensibilidade e uma avaliação experimental
contra doze meta-heurísticas conhecidas. Por fim, testam o método em quatro
problemas de projeto de engenharia e disponibilizam publicamente sua
implementação em MATLAB (Trojovská; Dehghani, 2022).

\section{Meta-heurística proposta}

Inicialmente, o CBOA gera aleatoriamente uma matriz com $N$ soluções, cada uma
formada por $m$ variáveis dentro dos respectivos limites. Após calcular a
função objetivo, as soluções são ordenadas: as $N_C$ melhores tornam-se chefs
instrutores e as demais tornam-se estudantes. O melhor indivíduo da população
é denominado melhor chef.

Na primeira fase, cada instrutor tenta melhorar sua solução de duas maneiras.
Na estratégia de exploração, ele se desloca em função da posição do melhor
chef, com perturbações aleatórias. Na estratégia de aproveitamento, realiza uma
busca local ao redor da própria posição. Os limites dessa vizinhança diminuem
com o número da iteração, produzindo passos progressivamente menores. Uma nova
posição somente substitui a anterior se melhorar o valor da função objetivo.

Concluída a atualização dos instrutores, o processo avança para a segunda
fase, dedicada aos estudantes, que escolhem aleatoriamente um instrutor e
atualizam sua posição seguindo três estratégias. Primeiro, aproxima todas as
suas variáveis das habilidades do instrutor selecionado. Depois, copia apenas
uma variável, interpretada como o aprendizado integral de uma habilidade. Por
fim, realiza uma pequena busca local em uma coordenada escolhida
aleatoriamente. Também nesse grupo prevalece uma regra gulosa: apenas
mudanças que melhoram a solução são aceitas. Ao final da iteração, toda a
população é novamente classificada e os papéis de instrutor e estudante podem
mudar. O processo é repetido por um número fixo de iterações. A complexidade
total do algoritmo depende do tamanho da população, do número de instrutores
e do número de iterações, sendo apresentada pelos autores em notação O grande
(Trojovská; Dehghani, 2022).

\section{Resultados obtidos}

O estudo empregou 23 funções clássicas: sete unimodais para avaliar
aproveitamento, seis multimodais de alta dimensão para avaliar exploração e dez
multimodais de dimensão fixa para avaliar o equilíbrio entre ambas. O CBOA e
os concorrentes foram executados vinte vezes, com mil iterações por execução.
Os indicadores foram média, melhor valor, desvio-padrão, mediana, tempo de
execução e posição no ranking. As comparações incluíram GA, PSO, GSA, TLBO,
GWO, MVO, WOA, TSA, MPA, HBA, DE e CMA.

Nos testes unimodais, o CBOA alcançou o ótimo global em cinco das sete funções
e obteve a primeira colocação geral. Nas funções multimodais de alta dimensão,
atingiu exatamente o ótimo em F9 e F11 e voltou a apresentar o melhor ranking
geral. Nas funções multimodais de dimensão fixa, também ficou em primeiro
lugar pela média das posições. O teste de soma de postos de Wilcoxon, com nível
de significância de 5\%, produziu valores de $p$ inferiores a 0,05 contra os
doze concorrentes, resultado interpretado pelos autores como superioridade
estatisticamente significativa (Wilcoxon, 1945; Trojovská; Dehghani, 2022).

O algoritmo também foi avaliado nas 29 funções da competição CEC 2017 (Awad et
al., 2016), obtendo a primeira colocação pelo ranking médio. Na análise de
sensibilidade, populações de 20, 30, 50 e 100 indivíduos produziram resultados
semelhantes em muitas funções, enquanto o aumento de 200 para mil iterações
geralmente melhorou as soluções. Nas aplicações reais --- projeto de vaso de
pressão, redutor de velocidade, viga soldada e mola de tração/compressão --- o
CBOA recebeu a primeira posição entre os algoritmos comparados nos quatro
casos.

\section{Conclusões}

Os autores concluem que a combinação das cinco estratégias permite ao CBOA
equilibrar a exploração de novas regiões com o refinamento local das soluções.
Os experimentos indicam desempenho competitivo em funções com características
distintas e em problemas restritos de engenharia. Entretanto, por ser um
método estocástico, o CBOA não garante o ótimo global e pode falhar em classes
de problemas não contempladas nos testes. Como trabalhos futuros, o artigo
sugere desenvolver versões binária e multiobjetivo e aplicar o algoritmo em
outros domínios. Assim, a principal conclusão é que o CBOA constitui uma
alternativa geral e relativamente simples para otimização contínua, embora sua
eficácia deva ser verificada experimentalmente para cada nova aplicação
(Trojovská; Dehghani, 2022).

\clearpage
\section*{REFERÊNCIAS}
\addcontentsline{toc}{section}{REFERÊNCIAS}

\begin{singlespace}
\noindent AWAD, N. H.; ALI, M. Z.; LIANG, J. J.; QU, B. Y.; SUGANTHAN, P. N.
\textbf{Evaluation criteria for the CEC 2017 special session and competition on
single objective real-parameter numerical optimization}. Daegu: Kyungpook
National University, 2016.

\vspace{\baselineskip}

\noindent TROJOVSKÁ, Eva; DEHGHANI, Mohammad. A new human-based metaheuristic
optimization method based on mimicking cooking training. \textbf{Scientific
Reports}, [s. l.], v. 12, art. 14861, 2022. DOI:
\href{https://doi.org/10.1038/s41598-022-19313-2}{10.1038/s41598-022-19313-2}.

\vspace{\baselineskip}

\noindent WILCOXON, Frank. Individual comparisons by ranking methods.
\textbf{Biometrics Bulletin}, [s. l.], v. 1, n. 6, p. 80--83, 1945.

\vspace{\baselineskip}

\noindent WOLPERT, David H.; MACREADY, William G. No free lunch theorems for
optimization. \textbf{IEEE Transactions on Evolutionary Computation}, [s. l.],
v. 1, n. 1, p. 67--82, 1997.
\end{singlespace}

\end{document}
